\chapter{Manage separations and spacing}\label{sec:separations}

The previous chapter covered cell and column management. We will now deal with how our cells and columns are separated. This chapter will be brief, enjoy it while it lasts \emoji{grinning}.

\section{Add all separation lines}
If you want to display all horizontal and/or vertical lines, there is a command for that. Since the option is passed to tabularray's header, do not forget to put the colspec inside the keyword argument \snippet{colspec=\{...\}}.

\begin{Code}[show-output]
\begin{tblr}{
	\highLight{hlines},  % displays all horizontal lines
	\highLight{vlines},  % displays all vertical lines
	colspec = {cccccccc},
}
	& Monday   & Tuesday  & Wednesday & Thursaday & Friday   & Saturday & Sunday \\
	Lunch  & Potatoes & Potatoes & Potatoes  & Potatoes  & Potatoes & Potatoes & Potatoes too! \\ 
	Dinner & Rice     & Stir-fry & Soup      & Eat out   & Fish     & Pizza    & Leftovers \\
\end{tblr}
\end{Code}

\section{Manage row and column spacing}
As we previously saw, tabularray rows are naturally larger than native Latex rows. This vertical spacing (as well as the horizontal spacing) can be changed with the options \snippet{rowsep} and \snippet{colsep} in the table's header. You can use any size indicator, it is common to use \snippet{pt} but \snippet{cm} or anything else that suits you also works.

\begin{Code}[show-output]
% reference table
\begin{tblr}{
	colspec={lll},
	hlines,
}
	cell 1 & cell 2             & cell $\frac{3}{3}$ \\
	cell 4 & cell $\frac{5}{6}$ & cell 6 \\
\end{tblr}
\end{Code}
\begin{Code}[show-output]
\begin{tblr}{
	colspec = {lll},
	hlines,
	\highLight{rowsep = {8pt},}  % default tabularray value is 2pt
	\highLight{colsep = {1pt},}  % default tabularray value is 6pt
}
	cell 1 & cell 2             & cell $\frac{3}{3}$ \\
cell 4 & cell $\frac{5}{6}$ & cell 6 \\
\end{tblr}
\end{Code}

As we can see, \snippet{rowsep} defines the margin between each row and its neighbours, both above and below. Similarly, \snippet{colsep} defines the left and right margin of columns.

\section{Manage merged cell expansion}

Let's have another look to merged cells. If you played with them a bit, you might have noticed that tabularray will expand your table so that the text in the merged cell fits on a single line. For instance:

\begin{Code}[show-output]
% Reference table
\begin{tblr}{
	colspec={llll},
	vlines,
	hlines,
}
	\SetCell[r=2]{m} {Mur, ville, \\ Et port, \\ Asile \\ De mort, \\ Mer grise \\ Où brise \\ La brise, \\Tout dort.} & Les Djinns & Victor Hugo &  1829 \\
	&\SetCell[c=3]{l} A beautiful poem & & \\
\end{tblr}
\end{Code}
\begin{Code}[show-output]
\begin{tblr}{
	colspec={llll},
	vlines,
	hlines,
}
	\SetCell[r=2]{m} {Mur, ville, \\ Et port, \\ Asile \\ De mort, \\ Mer grise \\ Où brise \\ La brise, \\Tout dort.} & Les Djinns & Victor Hugo &  1829 \\
	&\SetCell[c=3]{l} A beautiful poem in which each stanza has a different metrical structure & & \\
\end{tblr}
\end{Code}

We can see that only the last column is extended to make the merged cell fit. Similarly, only the last row has been extended to make the first cell fit. 

There are several ways to fix the issue. The first one is to set the columns (and rows) width, as merged cells respect their parent's column/row width, but this is not always interesting or easy. The second solution is to use another computation mode for cell expansion, with the header attribute \snippet{hspan} (and \snippet{vspan}).

\begin{itemize}
	\item \snippet{hspan=default} will extend the last column as much as possible, it is the default behaviour we saw.
	\item \snippet{hspan=even} will evenly distribute the extra space between the margins of the columns in which the cells are merged
	\item \snippet{hspace=minimal} will force the merged cell to not extend, and thus insert line breaks in the merged cell if necessary
\end{itemize}

Similarly, \snippet{vspan} can be either \snippet{default} or \snippet{even} (but not \snippet{minimal}!).

Some examples:

\begin{Code}[show-output]
\begin{tblr}{
		colspec={llll},
		vlines,
		hlines,
		\highLight{hspan=even},
		\highLight{vspan=even},
	}
	\SetCell[r=2]{m} {Mur, ville, \\ Et port, \\ Asile \\ De mort, \\ Mer grise \\ Où brise \\ La brise, \\Tout dort.} & Les Djinns & Victor Hugo &  1829 \\
	&\SetCell[c=3]{l} A beautiful poem in which each stanza has a different metrical structure & & \\
\end{tblr}
\end{Code}

\begin{Code}[show-output]
\begin{tblr}{
		colspec={llll},
		vlines,
		hlines,
		\highLight{hspan=minimal},
		\highLight{vspan=even},
	}
	\SetCell[r=2]{m} {Mur, ville, \\ Et port, \\ Asile \\ De mort, \\ Mer grise \\ Où brise \\ La brise, \\Tout dort.} & Les Djinns & Victor Hugo &  1829 \\
	&\SetCell[c=3]{l} A beautiful poem in which each stanza has a different metrical structure & & \\
\end{tblr}
\end{Code}

Note that \snippet{hspan=even} (or \snippet{vspan}) does not guarantee that the columns (or rows) will have the same width or height. It guarantees that the  \emph{margins} of these columns (or rows) will be the same. For instance, in the last table above, you can see that the two rows do not have the same height, even despite the \snippet{vspan=even}. However, if you measure the space between the text and the above/below borders, you will see it is the same in the first and second row.

\begin{infobox}{TL;DR}
\begin{itemize}
	\item Use \snippet{hlines}, \snippet{vlines} in the table header to automatically display horizontal and vertical lines
	\item Use \snippet{rowsep}, \snippet{colsep} in the table header to change the spacing between rows and columns
	\item Use \snippet{hspan} and \snippet{vspan} with the values \snippet{default}, \snippet{even} or (for \snippet{hspan} only) \snippet{minimal} to manage the extension of rows and columns containing merged cells.
\end{itemize}

\textbf{\underline{Exercise:}} recreate the following table.

\begin{center}
\begin{standalonebox}
\begin{tblr}{
		hlines,
		vlines,
		colsep=8pt,
		rowsep=6pt,
		hspan=minimal,
		vspan=even,
		colspec={Q[l,2cm,m]lll}
	}
	Title & Author & Year & \SetCell[r=2]{} Excerpt \\
	\SetCell[c=3]{c} Comment & & & \\
	Les Djinns & Victor Hugo &  1829 & \SetCell[r=2]{m} {Mur, ville,\\ Et port, \\ Asile \\ De mort, \\ Mer grise \\ Où brise \\ La brise, \\Tout dort.} \\
	\SetCell[c=3]{l} A beautiful poem in which each stanza has a different metrical structure. & & & \\
	Colloque sentimental & Paul Verlaine & 1869 & \SetCell[r=2]{m} {Dans le vieux parc solitaire et glacé \\ Deux formes ont tout à l'heure passé. \\[1em] Leurs yeux sont morts et leurs lèvres sont molles, \\ Et l'on entend à peine leurs paroles.} \\
	\SetCell[c=3]{l} A conclusion to the compilation ``Fêtes galantes'' that stands in stark contrast by its darkness and melancholy. \\
\end{tblr}
\end{standalonebox}
\end{center}

The columns must have a \snippet{8pt} separation, rows a \snippet{6pt} separation, the first column must be \snippet{2cm} wide. To have two consecutive line breaks in a single cell, you can use the command \snippet{\textbackslash\textbackslash[1em]}.

Solution is available in \Cref{app:separations_solution}
\end{infobox}